% part: intuitionistic-logic
% chapter: propositions-as-types
% section: proofs-to-terms

\documentclass[../../../include/open-logic-section]{subfiles}

\begin{document}

\olfileid{int}{pty}{pt}

\olsection{تحويل \usetoken{P}{derivation} إلى حدود برهان}

عملية تحويل~!!{proof}s الاستنباط الطبيعي إلى حدود برهان، أي ترجمة
!!{proof}s في~\Log{N2i} إلى~!!{proof}s في~\Log{tN2i}، مباشرة للغاية. فهي
لا تعدو كتابة حد برهان إلى يسار كل~!!{formula} في يمين متتالية ضمن
!!{derivation}، فتنتج متتاليات من الصورة~$\Gamma\Sequent M:!A$. ونقول عندئذ
إن~$M$ \emph{يشهد} لـ~$!A$ في السياق~$\Gamma$. وتحدد قواعد~\Log{tN2i}
تمامًا أي حد برهان نكتب.

\begin{prop}
كل~!!{proof}~$\delta$ للمتتالية~$\Gamma\Sequent !A$ في~\Log{N2i} يقابل
!!a{proof}~$\delta'$ للمتتالية~$\Gamma\Sequent N:!A$ في~\Log{tN2i}.
ويتحدد حد البرهان~$N$ تحديدًا وحيدًا بـ~$\delta$.
\end{prop}

\begin{proof}
نستعمل الاستقراء في ارتفاع~$\delta$. فإذا كان~$\delta$ المسلّمة
$x:!A\Sequent !A$، كان~$\delta'$ المسلّمة~$x:!A\Sequent x:!A$. وإذا انتهى
$\delta$ بقاعدة استدلال، وفر فرض الاستقراء حدود برهان لكل مقدمة، و!!{proof}s
في~\Log{tN2i} لمتتاليات تحتوي حدود البرهان هذه وتقابل كل مقدمة. ونطبق
القاعدة المقابلة في~\Log{tN2i}؛ فتحدد قاعدة استدلال~\Log{tN2i} حد البرهان
المقابل.
\end{proof}

\begin{ex}
  انظر في هذا~!!{proof} البسيط جدًا في~\Log{N1i} لـ
  $(!A\land !B)\lif(!B\land !A)$:
  \[
  \AxiomC{$\Discharge{!A\land !B}{x}$}
  \RightLabel{\Elim\land[2]}
  \UnaryInfC{$!B$}
  \AxiomC{$\Discharge{!A\land !B}{x}$}
  \RightLabel{\Elim\land[1]}
  \UnaryInfC{$!A$}
  \RightLabel{\Intro\land}
  \BinaryInfC{$!B \land !A$}
  \DischargeRule{\Intro\lif}{x}
  \UnaryInfC{$(!A \land !B) \lif (!B \land !A)$}
  \DisplayProof
  \]
  ويبدو في~\Log{N2i} هكذا:
  \[
  \Axiom$x : !A\land !B \fCenter !A\land !B$
  \RightLabel{\Elim\land[2]}
  \UnaryInf$x : !A\land !B \fCenter !B$
  \Axiom$x : !A\land !B \fCenter !A\land !B$
  \RightLabel{\Elim\land[1]}
  \UnaryInf$x : !A\land !B \fCenter !A$
  \RightLabel{\Intro\land}
  \BinaryInf$x : !A\land !B \fCenter !B \land !A$
  \RightLabel{\Intro\lif}
  \UnaryInf$\fCenter (!A \land !B) \lif (!B \land !A)$
  \DisplayProof
  \]
  نسند حدود البرهان من الأعلى إلى الأسفل. فحدا برهان المسلّمتين هما نفس
  المتغير في السياق، أي~$x$. ثم يتحدد حدا البرهان المرتبطان بـ
  $\Elim\land[i]$ على أنهما~$\proj{2}{x}$ و$\proj{1}{x}$. وعند تطبيق
  \Intro{\land} نجمع هذين الحدين:
  $\pair{\proj{2}{x}}{\proj{1}{x}}$. وأخيرًا تطبق قاعدة~\Intro{\lif}
  تجريد~$\lambd$، فتقيّد~$x$ وتسجل أن~$x$ وسم الفرض~$!A\land !B$:
  $\lambd[\typeof{x}{!A \land
  !B}][\pair{\proj{2}{x}}{\proj{1}{x}}]$.
  وعندئذ يكون~!!{proof} في~\Log{tN2i} هو:
  \[
  \Axiom$x : !A\land !B \fCenter x: !A\land !B$
  \RightLabel{\Elim\land[2]}
  \UnaryInf$x : !A\land !B \fCenter \proj{2}{x} : !B$
  \Axiom$x : !A\land !B \fCenter x: !A\land !B$
  \RightLabel{\Elim\land[1]}
  \UnaryInf$x : !A\land !B \fCenter \proj{1}{x} : !A$
  \RightLabel{\Intro\land}
  \BinaryInf$x : !A\land !B \fCenter \pair{\proj{2}{x}}{\proj{1}{x}} : 
    !B \land !A$
  \RightLabel{\Intro\lif}
  \UnaryInf$\fCenter \lambd[\typeof{x}{!A \land
  !B}][\pair{\proj{2}{x}}{\proj{1}{x}}] : (!A \land !B) \lif (!B \land !A)$
  \DisplayProof
  \]
\end{ex}


\end{document}
